% Chapter 29: The Geometry of Alphabets
\chapter{The Geometry of Alphabets}

\chapterepigraph{%
  An alphabet is not a list.\\
  It is a tiling of perceptual space\\
  whose pieces must be large enough to distinguish\\
  and small enough to learn.
}{Flyxion, \textit{Semantic Infrastructure}}

\sidenote{Connects to\\constraint basins\\(Ch.\,24) and\\the Boolean\\lattice (App.\,D)}

A writing system imposes a geometric structure on the space of possible
marks. The letters of an alphabet are not arbitrary — they are the result
of an optimization process operating over centuries, selecting for marks
that are perceptually distinct, motorically accessible, and spatially
efficient.

This chapter analyzes that structure using the accessibility and
constraint frameworks developed in earlier parts.

\section{Letters as Compressed Shapes}

Every letter is a compression: it discards the continuous variation of
handwriting and retains only the features necessary for identification.
But what features are retained?

\begin{definition}{Distinguishing Features of a Glyph}{glyph-features}
  The \emph{distinguishing feature set} of a glyph $g$ is the minimal
  set $F(g)$ of geometric properties such that $g$ is the unique glyph
  in the alphabet with all properties in $F(g)$:
  \[
    F(g) = \arg\min_{F' \subseteq \mathcal{F}} |F'| \text{ s.t. }
    \forall g' \neq g:\, F(g') \not\supseteq F'
  \]
  where $\mathcal{F}$ is the full space of geometric features.
\end{definition}

Different writing systems use different feature sets. The Latin alphabet
relies heavily on: vertical strokes (I, l), horizontal strokes (E, F, H),
diagonal strokes (A, V, W, X), curves (C, G, S, O), and enclosed regions
(B, D, O, P, R). Arabic relies on: dot count, dot position, baseline
connectivity, and curve direction. Chinese relies on: stroke count,
stroke type, and spatial arrangement.

\section{Symmetry, Rotation, Reflection}

Many alphabets exploit or avoid visual symmetry systematically.

\begin{proposition}{Symmetry Creates Confusion}{sym-confusion}
  Glyphs that are related by rotation or reflection are easily confused.
  If $g$ and $\rho(g)$ (where $\rho$ is a rotation or reflection) are
  both in the alphabet, their constraint basins $\mathcal{B}(g)$ and
  $\mathcal{B}(\rho(g))$ are adjacent in trace space, increasing
  confusion probability.
\end{proposition}

This explains a persistent difficulty in reading acquisition: children
confuse $b$ and $d$ (reflection), $p$ and $q$ (reflection), $b$ and $p$
(rotation), $n$ and $u$ (rotation). These pairs are related by simple
geometric transformations. A better-designed alphabet would maximize the
distance between all pairs of letters under all natural transformations.

The Arabic alphabet largely avoids this problem. Letters that share a
base shape are distinguished by diacritics (dot count and position)
rather than reflection or rotation, making the distinguishing feature
more salient and less prone to transformation confusion.

\section{Stroke Economy}

Every additional stroke in a glyph increases the motor acquisition
cost and the writing speed cost. The most economical alphabet uses
the minimum number of strokes to maintain perceptual distinctiveness.

\begin{definition}{Stroke Economy}{stroke-economy}
  The \emph{stroke economy} of an alphabet $\mathcal{A}$ is:
  \[
    E(\mathcal{A}) = \frac{|\mathcal{A}|}{
      \sum_{g \in \mathcal{A}} \text{strokes}(g)}
  \]
  — the number of letters per average stroke count. Higher economy
  means more letters per stroke.
\end{definition}

The Latin alphabet has moderate stroke economy: most letters require
2--4 strokes. Korean hangul has high stroke economy through its
systematic block structure: basic strokes are combined in consistent
patterns. Chinese characters have low stroke economy on average
(many strokes per character) but high semantic economy (each character
carries more meaning).

\section{The Accessibility Landscape of Alphabets}

An alphabet defines an accessibility landscape over the space of
possible marks. Each letter is an attractor; its basin is the
region of marks that readers will classify as that letter.

A well-designed alphabet maximizes the minimum basin size (all letters
are equally distinguishable) while minimizing the total feature complexity
(letters are as simple as possible). This is a classical information-theoretic
optimization, but with a geometric formulation.

\begin{namedthm}{The Alphabet Optimization Principle}
  An optimal alphabet $\mathcal{A}^*$ for a community with perceptual
  variability $\sigma$ and motor precision $\mu$ maximizes:
  \[
    \mathcal{A}^* = \arg\max_{\mathcal{A}} \min_{g \neq g' \in \mathcal{A}}
    d_\text{perc}(g, g') - \lambda \sum_{g \in \mathcal{A}} \text{complexity}(g)
  \]
  where $d_\text{perc}$ is perceptual distance and $\lambda$ trades
  distinctiveness against simplicity. Observed alphabets approximate
  this optimum, with the approximation improving over centuries of use.
\end{namedthm}

\begin{exercise}
  For the ten most common English letters (E, T, A, O, I, N, S, H, R, D),
  compute the pairwise Hamming distances in a simple feature space
  (presence/absence of: vertical stroke, horizontal stroke, diagonal,
  curve, enclosure, dot). Which pairs have minimum distance? Do these
  predict the most common reading confusions?
\end{exercise}

\begin{exercise}
  Design a minimal 26-letter alphabet that maximizes stroke economy
  while maintaining perceptual distinctiveness. How many strokes per
  letter does your alphabet average? Compare to the Latin alphabet.
\end{exercise}

\begin{exercise}[title={Cross-linguistic}]
  Compare the distinguishing feature sets of Arabic, Hebrew, Greek,
  and Latin scripts. Each developed from related Semitic origins.
  What geometric features did each preserve and which did each innovate?
  Can you identify a common ancestor feature set?
\end{exercise}
